\chapter{Комп'ютер із живих клітин}
% full version: chapters/21_living.tex

\pencilfig{21}{\pencilart{21}}{Живі нейрони на чипі: під ними тисячі електродів, і кожен уміє і слухати, і говорити.}

Чи можна побудувати машину не з кремнію, а з живих нейронів? Можна, і її вже будують. Нейрони вирощують просто на чипі з тисячами електродів --- пласким шаром або грудочками, схожими на крихітний мозок. Кожен електрод уміє і слухати, і подавати сигнал. Це стенд із відкритою схемою: усе можна виміряти й у все можна втрутитися. Тільки радіостанцій у ньому немає --- це шина даних без регістрів керування.

\section*{Що мережа робить сама}

Полишена сама на себе, культура час від часу спалахує цілком: майже всі нейрони клацають разом, а потім затихають, --- з проміжками від секунди до хвилин. Це вже знайомий генератор: мережа збуджує сама себе і втомлюється, доки з'єднання не відновлять запас речовини.

\pencilfig{129}{%
% spike raster of 8 neurons in a culture left to itself: sparse single clicks and shared bursts
\foreach \b in {0.9,3.1,4.0,6.6}{\fill[pcAmber, opacity=0.18] (\b-0.1,-0.15) rectangle (\b+0.32,2.95);}
\foreach \y/\ss in {0/{0.96,1.0,1.13,2.43,3.0,3.12,3.24,4.02,4.09,4.15,6.66,6.69,6.81},
  1/{0.46,0.88,1.0,1.05,3.09,3.2,3.94,4.04,4.37,6.65,6.71},
  2/{0.73,0.84,0.94,1.41,3.08,3.12,3.21,4.05,4.06,4.17,6.59,6.63},
  3/{0.89,0.93,1.06,3.1,3.19,3.28,3.89,3.94,4.13,4.2,4.31,6.63,6.65,6.78},
  4/{0.87,0.96,3.09,3.15,3.23,3.73,3.96,4.06,4.19,5.98,6.66,6.68,6.75},
  5/{0.44,0.92,0.97,3.05,3.22,3.27,4.01,4.07,4.17,5.76,6.57,6.73,6.75},
  6/{0.92,1.0,1.73,2.85,3.18,3.19,4.0,4.07,6.53,6.66,6.73},
  7/{0.87,0.92,3.13,3.13,3.27,3.99,4.06,4.17,6.44,6.61,6.72,7.13}}{
  \foreach \s in \ss {\draw[pcBlue, line width=0.7pt] (\s,0.35*\y) -- (\s,0.35*\y+0.25);}}
\draw[arr] (0,-0.3) -- (7.5,-0.3); \node[lbl, anchor=north] at (3.75,-0.32) {час};
\node[lbl, anchor=east, align=right] at (-0.05,1.35) {вісім\\нейронів};
\node[lbl, text=pcAmber!70!black, anchor=south] at (3.55,2.95) {спільні спалахи};
}{Культура на чипі, полишена сама на себе. Поодинокі клацання рідкісні, а час від часу майже вся мережа спалахує разом --- і затихає, доки з'єднання не відновлять запас.}

\section*{Учитель без дофаміну}

Як навчити таку мережу, якщо дофаміну немає? Перші способи були електричними. Наприклад: стимулювати мережу, доки вона не видасть потрібної відповіді, --- і одразу зупинитися. Відповідь закріплюється. В іншому досліді культура керувала простим віртуальним «тілом» і вчилася вести його до мети. У всіх цих дослідах учитель --- малюнок струму, який обирає інженер.

\section*{Дофамін за спалахом світла}

Одного разу вчителя спробували зробити хімічним. На одній із платформ із живими грудочками нейронів у поживну рідину додали дофамін у світлочутливій «клітці», яка не дає йому діяти. Коли мережа відповідала на стимул сильніше, ніж раніше, її «винагороджували»: спалах ультрафіолету руйнував клітку, і дофамін виходив назовні. За двісті сорок повторів відповіді загалом зростали. Але самі автори чесно пишуть, що з одного такого досліду не можна сказати, чи причина саме в дофаміні.

\section*{Жердина на візку}

У недавньому досліді грудочки нейронів навчали балансувати жердиною на візку --- класична задача для програм навчання. Які навчальні імпульси подавати, обирала програма. Сигнали, обрані програмою, працювали краще за випадкові, але після трьох чвертей години відпочинку вивчене зникало. А якщо заблокувати \nt{GLUT}-з'єднання, навчання не йшло зовсім: вивчене живе саме там. Учитель же, як і раніше, ззовні --- тільки тепер це не інженер, а програма.

\section*{Живий резервуар}

Є й інший підхід: живу мережу не навчають зовсім. Її використовують як «резервуар» --- багату, заплутану систему, чиї відповіді простий електронний лічильник учиться розбирати. Так жива грудочка допомагала розрізняти голоси. Щоправда, за словами авторів, зв'язки в самій грудочці під час роботи теж змінювалися.

\section*{Ціна і питання}

Мільйон нейронів на чипі споживає приблизно десятитисячну частку вата --- значно менше, ніж електроніка навколо них. І стенд не вміє сам вирішити, що вважати успіхом: цю роль завжди грає хтось ззовні. А в міру того як такі системи ростуть, постає й питання етики, якого поки ніхто не розв'язав.

\fullversion{21}
